\documentclass[11pt,a4paper]{article}
\usepackage[margin=24mm]{geometry}
\usepackage[hidelinks]{hyperref}
\title{Snow cover project: experiment record}
\author{Ahmad Ammar}
\date{Internal. Newest last.}
\begin{document}
\maketitle
Each entry gives the question, what changed, what it produced, how it was judged, and what it led to. The public
manuscript reports only the final method and its benchmarks; the reasoning lives here. Numbers come from the
generated values of the report (\texttt{docs/generated}), so this record and the report cannot drift apart.

\section{The original series (MS-era work)}
\textbf{What existed.} One summary table, \texttt{UIBNDSI.xlsx}: 909 records of snow area from 26 February 2000 to
19 December 2019, from 8-day MODIS composites. The processing before the table was not kept, and the project's documents
disagree about the boundary of the Upper Indus. See \texttt{original\_audit.tex}.
\textbf{Judgement.} The table can be compared with, but not rebuilt into, a reproducible result. Statements in the old
report about years outside the series cannot be supported by it.
\textbf{Led to.} A rebuild from source data with a documented boundary.

\section{Rebuild: 2001--2025, daily data, per-pixel change}
\textbf{Change.} Daily MOD10A1 collection 6.1 at 500~m, hydrological years October to September, 2001--2025. Regions
from GMBA mountain divisions and HydroBASINS sub-basins, replacing rectangles. Each pixel with at least 20 usable years
tested with Mann-Kendall, with false discovery control at 5\%.
\textbf{Why.} The 8-day product hides cloud gaps, and a regional mean hides opposite changes within a region. Per-pixel
classes (gain, loss, no detectable trend, persistent, rarely snow, too little data) show where change happens.
\textbf{Corrections on the way.} The Hamed-Rao correction gave near-zero $p$ for very weak trends in 25 values, and was
dropped. A coarse histogram made the false discovery threshold return zero significant pixels wrongly; a finer one
fixed it, and all ranges were rerun.
\textbf{Result.} After correction, loss is concentrated in the Himalaya beside the cloudiest ground, little area gains,
and almost no range or sub-basin mean shows a detectable trend. Loss is robust to using Aqua instead of Terra and to the
snow threshold tested so far.

\section{The literature, and what the question becomes}
\textbf{Change.} Scopus queries and a citation corpus (\texttt{references/}), read against the provisional findings.
\textbf{Result.} Regional trend studies exist and disagree. See \texttt{gap\_analysis.tex}.
\textbf{Led to.} The main question moves from ``is snow changing'' to why published trends disagree, and how large a
change can be detected with each sensor and period.

\section{Further back in time}
\textbf{Question.} Can the record start before 2000?
\textbf{Change.} Landsat availability checked back to the 1970s; MSS has no shortwave infrared band, so it cannot
compute NDSI in the same way. Harmonisation between Landsat sensors fitted from near-same-day pairs over the region,
compared with the published coefficients of Roy et al.\ (2016).
\textbf{Status.} Landsat can extend selected patches, not the whole region, before 2000. Still to do: a Landsat check on
the Himalayan loss patches, and the AVHRR overlap.

\section{Open steps}
Detectability analysis; comparison with MOYDGL06; reasons published trends disagree, as a set of tested explanations.
\end{document}
